\section{Los datos}\label{sec:datos}
Los organizadores entregaron cinco fuentes de información (Tabla~\ref{tab:fuentes}).

\begin{table*}[!t]
\caption{Fuentes de información entregadas}\label{tab:fuentes}
\centering\footnotesize
\begin{tabularx}{\textwidth}{@{}>{\raggedright\arraybackslash}p{3.2cm}X>{\raggedright\arraybackslash}p{2.6cm}>{\raggedright\arraybackslash}p{2.4cm}@{}}
\toprule
Fuente & Contenido & Resolución & Periodo \\
\midrule
Parcelas & 197 polígonos con identificador, área, municipio, estado y conjunto & Exacta & Ciclo 2025 \\
Conjunto Básico & 23 índices de vegetación por parcela y fecha, con media, desviación estándar, máximo y mínimo dentro de la parcela, de Sentinel-2 y Landsat & 10 a 30 m, cada 5 a 16 días & 2022 a 2025 \\
Conjunto PRO & NDVI, EVI, LAI y MSAVI por parcela y fecha, de Planet & Unos 3 m, casi diario & 2025 completo \\
Clima & Precipitación mensual (CHIRPS~\cite{funk2015}) y temperatura mínima y máxima mensual (CHIRTS-ERA5~\cite{chirts}) & Unos 5 km & 2022 a 2025 \\
Topografía & Elevación y pendiente del Continuo de Elevaciones Mexicano~\cite{inegi} & 120 m & Estática \\
\bottomrule
\end{tabularx}
\end{table*}

\textbf{Parcelas.} Las 197 parcelas están en nueve municipios, dentro de un área de unos 42 por 50 kilómetros.
Chignahuapan, en Puebla, concentra 100. La parcela típica mide alrededor de 5 hectáreas y la más chica menos de
media hectárea.

\textbf{Imágenes satelitales.} Cada satélite mide con su propio sensor, así que sus índices no se deben mezclar sin
ajustarlos. Por ejemplo, Planet da valores de NDVI más altos que Sentinel-2 sobre el mismo suelo. El campo de
nubosidad vale casi siempre 0 o 100, y cuando vale 100 los índices vienen vacíos. Aun así, la cobertura es buena.
Entre mayo y octubre de 2025 cada parcela tiene entre 19 y 34 observaciones limpias de Sentinel-2, más que en
cualquiera de los tres años anteriores. Además, el LAI de Planet se calcula con su banda de borde rojo
(\emph{red-edge}), así que de él se puede recuperar el índice NDRE, que la literatura señala como uno de los más
útiles para estimar rendimiento.

\textbf{Fechas.} Las fechas vienen en formato día, mes y año. Si se ordenan como texto y no como fechas, Planet
parece terminar el 31 de agosto, cuando en realidad llega al 30 de diciembre de 2025.

\textbf{Clima y topografía.} El clima es mensual, así que no permite medir lo que pasó en una semana concreta, como
la de la floración. Sus celdas miden unos 5 kilómetros y las 197 parcelas caen en sólo 26 celdas, por lo que el
clima funciona casi como una etiqueta de zona. La topografía se entregó a 120 metros, aunque la descripción oficial
habla de 15, y una parcela típica cubre apenas unas tres celdas.
